% ── chemistry-handout-preamble.tex ── 化学竞赛讲义 LaTeX 导言区（v7 · 精排版）
\documentclass[a4paper,11pt]{ctexart}

% ── 字体优化 ──
\setCJKmainfont[AutoFakeBold=true,BoldFont=SimHei]{SimSun}
\setCJKsansfont{SimHei}

% 化学式宏（mhchem）
\usepackage[version=4]{mhchem}

% ── 数学 ──
\usepackage{amsmath,amssymb}

% ── 图片 ──
\usepackage{graphicx}

% ── 表格（跨页+自动列宽）──
\usepackage{booktabs}
\usepackage{xltabular}
\usepackage{array}
\newcolumntype{L}{>{\raggedright\arraybackslash}X}
\newcolumntype{C}{>{\centering\arraybackslash}X}
\newcolumntype{R}{>{\raggedleft\arraybackslash}X}

% ── 浮动体（[H] 强制定位）──
\usepackage{float}

% ── 页面 ──
\usepackage[includeheadfoot,margin=2.2cm,headheight=14pt]{geometry}
\setlength{\textwidth}{16.6cm}

% ── 行间距 ──
\usepackage{setspace}
\setstretch{1.15}
\setlength{\parskip}{0.3em}
\setlength{\parindent}{2em}

% ── 防止表格跨页 ──
\usepackage{longtable}
\setlength{\LTpre}{0pt}
\setlength{\LTpost}{0pt}

% ── 防止图片跨页（尽量）──
\renewcommand{\textfraction}{0.1}
\renewcommand{\topfraction}{0.9}
\renewcommand{\bottomfraction}{0.9}
\setcounter{totalnumber}{10}
\setcounter{topnumber}{10}
\setcounter{bottomnumber}{10}

% ── 超链接(无红框) ──
\usepackage[colorlinks=true,linkcolor=black,citecolor=black,urlcolor=black]{hyperref}

% ── 颜色 ──
\usepackage[dvipsnames,svgnames]{xcolor}
\definecolor{chapterblue}{RGB}{23,50,77}
\definecolor{insightbg}{RGB}{235,240,250}
\definecolor{warningbg}{RGB}{255,245,235}
\definecolor{tipbg}{RGB}{235,250,235}

% ── 页眉页脚 ──
\usepackage{fancyhdr}
\pagestyle{fancy}
\fancyhf{}
\fancyhead[L]{\small\color{chapterblue}\leftmark}
\fancyhead[R]{\small\color{chapterblue}\thepage}
\renewcommand{\headrule}{{\color{chapterblue}\hrule height 0.8pt}}

% ── 章节编号与样式 ──
\setcounter{secnumdepth}{1}
\ctexset{
  section = {
    format = \Large\bfseries\heiti\color{chapterblue},
    aftername = \hspace{0.5em},
    beforeskip = 2ex,
    afterskip = 1ex,
  },
  subsection = {
    format = \large\bfseries\heiti,
    beforeskip = 1.5ex,
    afterskip = 0.5ex,
  },
  subsubsection = {
    format = \normalsize\bfseries\heiti,
    beforeskip = 1ex,
    afterskip = 0.3ex,
  },
  paragraph = {
    format = \small\bfseries,
    beforeskip = 0.5ex,
    afterskip = 0.2ex,
  },
}

% ── blockquote（引用/提示框）缩减 ──
\usepackage{etoolbox}
\AtBeginEnvironment{quote}{\small\setstretch{1.05}\leftskip=1em\rightskip=1em}

% ── 教学提示框（浅底色版）──
\newenvironment{insightbox}{%
  \par\smallskip
  \noindent{\small\bfseries\color{RoyalBlue}🧠 认知冲突}\par
  \leavevmode\colorbox{insightbg}{\parbox{\dimexpr\textwidth-2\fboxsep\relax}{%
  \ignorespaces}}%
}{\par\smallskip}

\newenvironment{warningbox}{%
  \par\smallskip
  \noindent{\small\bfseries\color{BrickRed}⚠️ 易错点}\par
  \leavevmode\colorbox{warningbg}{\parbox{\dimexpr\textwidth-2\fboxsep\relax}{%
  \ignorespaces}}%
}{\par\smallskip}

\newenvironment{tipbox}{%
  \par\smallskip
  \noindent{\small\bfseries\color{ForestGreen}💡 理解提示}\par
  \leavevmode\colorbox{tipbg}{\parbox{\dimexpr\textwidth-2\fboxsep\relax}{%
  \ignorespaces}}%
}{\par\smallskip}

% ── 列表（答案区加间距）──
\usepackage{enumitem}
\setlist{nosep}
% 专门用于参考答案的列表（每项之间空一行）
\newlist{answerlist}{enumerate}{1}
\setlist[answerlist]{label=\arabic*., leftmargin=2em, itemsep=0.8em}

% ── 防止内容溢出 ──
\usepackage{microtype}
\setlength{\emergencystretch}{2em}
\setlength{\hfuzz}{2pt}
\setlength{\overfullrule}{0pt}
\sloppy

% ── pandoc 兼容 ──
\providecommand{\tightlist}{\setlength{\itemsep}{0pt}\setlength{\parskip}{0pt}}
\newenvironment{Shaded}{}{}
\newenvironment{Highlighting}{}{}
\newcommand{\NormalTok}[1]{#1}
\newcommand{\CommentTok}[1]{#1}
\newcommand{\KeywordTok}[1]{#1}
\newcommand{\StringTok}[1]{#1}
\newcommand{\FunctionTok}[1]{#1}
\newcommand{\DataTypeTok}[1]{#1}
\newcommand{\NumberTok}[1]{#1}
\newcommand{\ControlFlowTok}[1]{#1}
\newcommand{\OperatorTok}[1]{#1}
\newcommand{\BaseNTok}[1]{#1}
\newcounter{none}

% ── 图片尺寸限制 ──
\AtBeginDocument{\setkeys{Gin}{width=\textwidth,keepaspectratio}}

% ── 防止图片跨节 ──
\usepackage{placeins}

% ── 自定义命令 ──
% 练习题标题（去掉括号）
% 例题标题（去掉中括号）
% 这些需要在convert.py预处理中做正则替换

\graphicspath{{C:/Obsidion/妙妙屋/11-模板/scripts/.chem_media/}}
\title{电化学基础}
\date{}
\setlength{\tabcolsep}{3.5pt}
\small
\begin{document}
\maketitle
\tableofcontents
\newpage
\section{电化学基础}\label{ux7535ux5316ux5b66ux57faux7840}

\textbf{本讲定位}：电化学是热力学的''电学翻译''---把 \(\Delta \1^{\circ}\) 翻译成可直接测量的 \(E^{\circ}\)，把 \(K^{\circ}\) 翻译成电池电动势。这是连接热力学、平衡和氧化还原的枢纽。
\textbf{上节衔接}：← 热力学初步-超级充实版(自学完整)\textbar 热力学初步(\(\Delta \1^{\circ}\) = −nFE° 的热力学基础)← 化学平衡-超级充实版(自学完整)\textbar 化学平衡(Q/K 比较方法)
\textbf{下节衔接}：→ 2026-06-23-酸碱理论\textbar 酸碱理论(酸碱平衡中的电化学应用)→ 2026-06-23-氧化还原反应\textbar 氧化还原反应(元素化学中的电化学应用)
\textbf{主框架教材}：《普通化学原理》第4版 第10章
\textbf{辅助教材}：上海中学竞赛课程第八讲、Atkins《物理化学》电化学章节、Zchem 物理化学下 §3.3

\begin{center}\rule{0.5\linewidth}{0.5pt}\end{center}

\subsection{学习目标}\label{ux5b66ux4e60ux76eeux6807}

\begin{itemize}
\tightlist
\item
  目标1：能写出四种常见电极类型的半反应，能用电池符号表示原电池
\item
  目标2：能查标准电极电势表比较氧化剂/还原剂强弱，用 \(E^{\circ}\) 判断反应方向
\item
  目标3：理解 Nernst 方程的推导链(van't Hoff 等温式 → ΔG=-nFE → E=\(E^{\circ}\)-(RT/nF)lnQ)
\item
  目标4：能用 Nernst 方程分析浓度、pH、沉淀、配合物对电极电势的影响
\item
  目标5：能用 \(\Delta \1^{\circ}\) = −nFE° 建立热力学与电化学的定量桥梁，理解 \(\Delta \1^{\circ}\)-\(E^{\circ}\)-\(K^{\circ}\) 三角关系
\item
  目标6：能读懂 Latimer 图，计算物种间的 \(E^{\circ}\)，判断歧化反应
\item
  目标7：了解电化学腐蚀的基本类型和防腐方法
\end{itemize}

\begin{center}\rule{0.5\linewidth}{0.5pt}\end{center}

\subsection{一、为什么需要电化学}\label{ux4e00ux4e3aux4ec0ux4e48ux9700ux8981ux7535ux5316ux5b66}

\textbf{动机段}：热力学告诉我们 \(\Delta \1^{\circ}\) = −RT ln \(K^{\circ}\)，但在实际化学中，很多反应的 \(\Delta \1^{\circ}\) 很难直接测量。电化学提供了一条捷径---把化学反应''拆''成两个半反应，让电子通过外电路流动，产生可直接测量的电压。例如：
- Zn + Cu\textsuperscript{2+} → Zn\textsuperscript{2+} + Cu 的 \(\Delta \1^{\circ}\) = −212 kJ/mol → \(E^{\circ}\) = 1.10 V(可直接用伏特计测量)
- \(E^{\circ}\) 的精度远高于量热法测 \(\Delta \1^{\circ}\)，因此电化学是获取热力学数据最精确的方法之一
反过来，已知 \(E^{\circ}\) 也能计算 \(K^{\circ}\)---这就是''电化学桥''的核心价值。

\textbf{氧化还原反应的本质}：有电子转移或偏移的反应。氧化态和还原态的共轭关系与酸碱共轭关系相似：

\[
\text{还原态} \rightleftharpoons \text{氧化态} + n\mathrm{e}^- \qquad \text{酸} \rightleftharpoons \text{碱} + n\mathrm{H}^+
\]

前者为电子的转移，后者为质子的转移。有失电子的一方，便有得电子的一方，电子的得失一定同时发生。

\begin{center}\rule{0.5\linewidth}{0.5pt}\end{center}

\subsection{二、原电池与电动势}\label{ux4e8cux539fux7535ux6c60ux4e0eux7535ux52a8ux52bf}

\subsubsection{2.1 原电池的基本结构}\label{ux539fux7535ux6c60ux7684ux57faux672cux7ed3ux6784}

原电池是将自发氧化还原反应的化学能转化为电能的装置。核心设计：\textbf{氧化反应和还原反应在空间上分离，电子通过外电路传输}。

{\def\LTcaptype{none} % do not increment counter
\begin{longtable}[]{@{}lll@{}}
\toprule\noalign{}
组件 & 功能 & Daniel 电池示例 \\
\midrule\noalign{}
\endhead
\bottomrule\noalign{}
\endlastfoot
\textbf{负极(阳极)} & 氧化反应，释放电子 & Zn(s) → Zn\textsuperscript{2+}(aq) + 2e\textsuperscript{-} \\
\textbf{正极(阴极)} & 还原反应，接收电子 & Cu\textsuperscript{2+}(aq) + 2e\textsuperscript{-} → Cu(s) \\
\textbf{外电路} & 电子传输路径 & 导线 \\
\textbf{盐桥} & 维持电荷平衡(离子迁移) & KNO\textsubscript{3} 琼脂 \\
\end{longtable}
}

\textbf{盐桥的作用}：没有盐桥，负极侧因 Zn\textsuperscript{2+} 积累而带正电，正极侧因 Cu\textsuperscript{2+} 消耗而带负电，电荷积累会阻碍电子流动，电池很快停止工作。盐桥中的 K\textsuperscript{+} 向正极迁移，NO\textsubscript{3}\textsuperscript{-} 向负极迁移，维持两半电池的电中性。

\subsubsection{2.2 电池符号表示}\label{ux7535ux6c60ux7b26ux53f7ux8868ux793a}

Daniel 电池的标准表示：

\[
\text{Zn(s)} \mid \text{Zn}^{2+}\text{(c}_1\text{)} \| \text{Cu}^{2+}\text{(c}_2\text{)} \mid \text{Cu(s)}
\]

\textbf{书写规则}：
- 负极写左边，正极写右边
- 单竖线''\textbar{}``表示相界面
- 双竖线''‖``表示盐桥
- 从左到右的物质顺序：还原态 \textbar{} 氧化态 ‖ 氧化态 \textbar{} 还原态

\subsubsection{2.3 四种常见电极类型}\label{ux56dbux79cdux5e38ux89c1ux7535ux6781ux7c7bux578b}

{\def\LTcaptype{none} % do not increment counter
\begin{longtable}[]{@{}
  >{\raggedright\arraybackslash}p{(\linewidth - 6\tabcolsep) * \real{0.2500}}
  >{\raggedright\arraybackslash}p{(\linewidth - 6\tabcolsep) * \real{0.2500}}
  >{\raggedright\arraybackslash}p{(\linewidth - 6\tabcolsep) * \real{0.2500}}
  >{\raggedright\arraybackslash}p{(\linewidth - 6\tabcolsep) * \real{0.2500}}@{}}
\toprule\noalign{}
\begin{minipage}[b]{\linewidth}\raggedright
电极类型
\end{minipage} & \begin{minipage}[b]{\linewidth}\raggedright
示例
\end{minipage} & \begin{minipage}[b]{\linewidth}\raggedright
半反应
\end{minipage} & \begin{minipage}[b]{\linewidth}\raggedright
电极符号
\end{minipage} \\
\midrule\noalign{}
\endhead
\bottomrule\noalign{}
\endlastfoot
\textbf{金属-金属离子} & Cu/Cu\textsuperscript{2+} & Cu\textsuperscript{2+} + 2e\textsuperscript{-} ⇌ Cu(s) & Cu(s) \textbar{} Cu\textsuperscript{2+}(aq) \\
\textbf{气体电极} & H\textsubscript{2}/H\textsuperscript{+} & 2H\textsuperscript{+} + 2e\textsuperscript{-} ⇌ H\textsubscript{2}(g) & Pt(s) \textbar{} H\textsubscript{2}(g) \textbar{} H\textsuperscript{+}(aq) \\
\textbf{难溶盐电极} & AgCl/Ag & AgCl(s) + e\textsuperscript{-} ⇌ Ag(s) + Cl\textsuperscript{-} & Ag(s) \textbar{} AgCl(s) \textbar{} Cl\textsuperscript{-}(aq) \\
\textbf{氧化还原电极} & Fe\textsuperscript{3+}/Fe\textsuperscript{2+} & Fe\textsuperscript{3+} + e\textsuperscript{-} ⇌ Fe\textsuperscript{2+} & Pt(s) \textbar{} Fe\textsuperscript{2+}(aq), Fe\textsuperscript{3+}(aq) \\
\end{longtable}
}

\textbf{气体电极和氧化还原电极需要惰性导体}(如 Pt)作为电子传输媒介，因为反应物和产物都是溶液中的离子或气体。

\subsubsection{2.4 电动势 E}\label{ux7535ux52a8ux52bf-e}

\[
E_{\text{cell}} = E_{\text{正极}} - E_{\text{负极}}
\]

当 \(E_{\text{cell}} > 0\) 时，反应自发。

\begin{center}\rule{0.5\linewidth}{0.5pt}\end{center}

\subsection{三、标准电极电势}\label{ux4e09ux6807ux51c6ux7535ux6781ux7535ux52bf}

\subsubsection{3.1 标准氢电极(SHE)}\label{ux6807ux51c6ux6c22ux7535ux6781she}

规定：在标准状态下，\(2\mathrm{H^+(aq, 1\ mol/L) + 2e^- \rightleftharpoons H_2(g, 100\ kPa)}\) 的电极电势 \(E^\theta = 0.000\ \mathrm{V}\)。

所有其他电极电势都是相对于 SHE 的\textbf{相对值}。

\subsubsection{3.2 标准电极电势表的使用}\label{ux6807ux51c6ux7535ux6781ux7535ux52bfux8868ux7684ux4f7fux7528}

\textbf{核心规则}：

{\def\LTcaptype{none} % do not increment counter
\begin{longtable}[]{@{}
  >{\raggedright\arraybackslash}p{(\linewidth - 2\tabcolsep) * \real{0.5000}}
  >{\raggedright\arraybackslash}p{(\linewidth - 2\tabcolsep) * \real{0.5000}}@{}}
\toprule\noalign{}
\begin{minipage}[b]{\linewidth}\raggedright
规则
\end{minipage} & \begin{minipage}[b]{\linewidth}\raggedright
说明
\end{minipage} \\
\midrule\noalign{}
\endhead
\bottomrule\noalign{}
\endlastfoot
\(E^{\circ}\) 越正 & 氧化态的氧化能力越强(越容易得电子) \\
\(E^{\circ}\) 越负 & 还原态的还原能力越强(越容易失电子) \\
\(E^{\circ}\) 是强度量 & 与半反应中电子数无关(乘以系数不改变 \(E^{\circ}\)) \\
正向 vs 逆向 & 半反应反向时 \(E^{\circ}\) 变号 \\
\end{longtable}
}

\textbf{判断反应自发性}：

\[
E_{\text{cell}}^\theta = E^\theta_{\text{氧化剂}} - E^\theta_{\text{还原剂}} > 0 \implies \text{反应自发}
\]

\begin{center}\rule{0.5\linewidth}{0.5pt}\end{center}

\subsection{四、Nernst 方程}\label{ux56dbnernst-ux65b9ux7a0b}

\subsubsection{4.1 完整推导链}\label{ux5b8cux6574ux63a8ux5bfcux94fe}

Nernst 方程可以从 van't Hoff 等温式严格推导：

\textbf{第1步}：van't Hoff 等温式

\[\Delta G = \Delta G^\circ + RT\ln Q\]

\textbf{第2步}：电化学功与自由能的关系

恒温恒压下，最大非体积功(电功)等于自由能变的减少：

\[\Delta G = W_{\text{电,max}} = -nEF\]

其中 n 为转移电子数，F = 96485 C/mol(Faraday 常数)。

\textbf{第3步}：代入等温式

\[-nFE = -nFE^\circ + RT\ln Q\]

\textbf{第4步}：整理得 Nernst 方程

\[\boxed{E = E^\circ - \frac{RT}{nF}\ln Q}\]

298 K 时简化为：

\[E = E^\theta - \frac{0.05916}{n}\lg Q\]

\subsubsection{4.2 Nernst 方程的典型应用}\label{nernst-ux65b9ux7a0bux7684ux5178ux578bux5e94ux7528}

\paragraph{(1) 浓度对电极电势的影响}\label{ux6d53ux5ea6ux5bf9ux7535ux6781ux7535ux52bfux7684ux5f71ux54cd}

\[
\mathrm{Zn^{2+} + 2e^- \rightleftharpoons Zn(s)} \quad E^\theta = -0.763\ \mathrm{V}
\]

当 \([\mathrm{Zn^{2+}}] = 0.001\) mol/L 时：

\[
E = -0.763 + \frac{0.05916}{2}\lg\frac{1}{0.001} = -0.763 + 0.089 = -0.674\ \mathrm{V}
\]

降低 {[}Zn\textsuperscript{2+}{]} 使 E 增大(更正)，即 Zn 的还原倾向减弱。

\paragraph{(2) pH 对电极电势的影响}\label{ph-ux5bf9ux7535ux6781ux7535ux52bfux7684ux5f71ux54cd}

\[
\mathrm{MnO_4^- + 8H^+ + 5e^- \rightleftharpoons Mn^{2+} + 4H_2O} \quad E^\theta = +1.51\ \mathrm{V}
\]

\[
E = 1.51 + \frac{0.05916}{5}\lg\frac{[\mathrm{MnO_4^-}][\mathrm{H^+}]^8}{[\mathrm{Mn^{2+}}]}
\]

pH 增大({[}H\textsuperscript{+}{]} 减小)→ E 显著降低 → KMnO\textsubscript{4} 氧化能力减弱。

\textbf{定量计算}：pH = 5 时(其他物种标准浓度)：

\[E = 1.51 + \frac{0.05916}{5}\lg(10^{-5})^8 = 1.51 - 0.473 = +1.037\ \mathrm{V}\]

pH 从 0 升到 5，E 降低 0.47 V。这就是为什么 KMnO\textsubscript{4} 在碱性溶液中氧化能力显著减弱。

\paragraph{(3) 沉淀对电极电势的影响}\label{ux6c89ux6dc0ux5bf9ux7535ux6781ux7535ux52bfux7684ux5f71ux54cd}

\[
\mathrm{Ag^+ + e^- \rightleftharpoons Ag(s)} \quad E^\theta = +0.799\ \mathrm{V}
\]

加入 Cl\textsuperscript{-} 生成 AgCl 沉淀后：

\[
E = 0.799 + 0.05916\lg\frac{K_{sp}}{[\mathrm{Cl^-}]}
\]

当 {[}Cl\textsuperscript{-}{]} = 1 mol/L 时 \(E = 0.799 + 0.05916 \times \lg(1.8 \times 10^{-10}) = +0.222\ \mathrm{V}\)(这就是 Ag/AgCl 参比电极的电势值)

\textbf{本质}：沉淀使 {[}Ag\textsuperscript{+}{]} 降低 → Nernst 方程预测 E 降低 → Ag 的还原能力减弱。

\paragraph{(4) 配合物对电极电势的影响}\label{ux914dux5408ux7269ux5bf9ux7535ux6781ux7535ux52bfux7684ux5f71ux54cd}

配合物形成使游离金属离子浓度降低 → E 降低。例如 {[}Cu(NH\textsubscript{3})\textsubscript{4}{]}\textsuperscript{2+}/Cu 的 \(E^{\circ}\) = −0.05 V(远低于 Cu\textsuperscript{2+}/Cu 的 +0.34 V)。

\textbf{实例计算}：Cu\textsuperscript{2+}/Cu 在 {[}Cu(NH\textsubscript{3})\textsubscript{4}\textsuperscript{2+}{]} = 0.1 mol/L、{[}NH\textsubscript{3}{]} = 1.0 mol/L 时：

游离 {[}Cu\textsuperscript{2+}{]} = {[}Cu(NH\textsubscript{3})\textsubscript{4}\textsuperscript{2+}{]}/(\(K_{\1}\) × {[}NH\textsubscript{3}{]}\textsuperscript{4}) = 0.1/(1.1×10\textsuperscript{13} × 1.0) = 9.1×10\textsuperscript{-15} mol/L

\[E = 0.34 + \frac{0.05916}{2}\lg(9.1\times10^{-15}) = 0.34 - 0.413 = -0.073\ \mathrm{V}\]

配合物形成使 Cu\textsuperscript{2+}/Cu 的电势从 +0.34 V 降至 −0.07 V。

\begin{center}\rule{0.5\linewidth}{0.5pt}\end{center}

\subsection{\texorpdfstring{五、\(\Delta \1^{\circ}\) = −nFE° 与 \(\Delta \1^{\circ}\) = −RT ln \(K^{\circ}\) 的统一}{五、\textbackslash Delta \textbackslash1\^{}\{\textbackslash circ\} = −nFE° 与 \textbackslash Delta \textbackslash1\^{}\{\textbackslash circ\} = −RT ln K\^{}\{\textbackslash circ\} 的统一}}\label{ux4e94delta-1circ-nfe-ux4e0e-delta-1circ-rt-ln-kcirc-ux7684ux7edfux4e00}

\subsubsection{5.1 三个核心桥梁公式}\label{ux4e09ux4e2aux6838ux5fc3ux6865ux6881ux516cux5f0f}

\[
\boxed{\Delta_rG^\circ = -nFE^\circ_{\text{cell}} = -RT\ln K^\theta}
\]

由此可得：

\[
E^\circ_{\text{cell}} = \frac{RT}{nF}\ln K^\theta = \frac{0.05916}{n}\lg K^\theta \quad (298\ \mathrm{K})
\]

{\def\LTcaptype{none} % do not increment counter
\begin{longtable}[]{@{}ccl@{}}
\toprule\noalign{}
\(E^{\circ}\)\_cell & \(K^{\circ}\) & 含义 \\
\midrule\noalign{}
\endhead
\bottomrule\noalign{}
\endlastfoot
\textgreater\textgreater{} 0 & \textgreater\textgreater{} 1 & 反应正向进行极其彻底 \\
≈ 0 & ≈ 1 & 反应物与产物共存 \\
\textless\textless{} 0 & \textless\textless{} 1 & 反应几乎不发生 \\
\end{longtable}
}

\subsubsection{\texorpdfstring{5.2 示例：从 \(E^{\circ}\) 计算 \(K^{\circ}\)}{5.2 示例：从 E\^{}\{\textbackslash circ\} 计算 K\^{}\{\textbackslash circ\}}}\label{ux793aux4f8bux4ece-ecirc-ux8ba1ux7b97-kcirc}

Daniel 电池：\(E^{\circ}\) = 1.10 V，n = 2

\[
\lg K^\theta = \frac{nE^\circ}{0.05916} = \frac{2 \times 1.10}{0.05916} = 37.2
\]

\[
K^\theta = 10^{37.2} \approx 1.6 \times 10^{37}
\]

\(K^{\circ}\) 极大，说明 Zn + Cu\textsuperscript{2+} → Zn\textsuperscript{2+} + Cu 几乎完全进行。

\subsubsection{\texorpdfstring{5.3 从 \(E^{\circ}\) 计算 \(K_{\1}\)}{5.3 从 E\^{}\{\textbackslash circ\} 计算 K\_\{\textbackslash1\}}}\label{ux4ece-ecirc-ux8ba1ux7b97-k_1}

\textbf{方法}：设计两个半反应，一个涉及难溶盐的溶解平衡，组合后求 \(K_{\1}\)。

例：已知 \(E^{\circ}\)(Ag\textsuperscript{+}/Ag) = +0.799 V，\(E^{\circ}\)(AgCl/Ag) = +0.222 V。

Ag\textsuperscript{+}/Ag：Ag\textsuperscript{+} + e\textsuperscript{-} → Ag(s)，\(E^{\circ}\) = +0.799 V
AgCl/Ag：AgCl(s) + e\textsuperscript{-} → Ag(s) + Cl\textsuperscript{-}，\(E^{\circ}\) = +0.222 V

两式相减：AgCl(s) → Ag\textsuperscript{+}(aq) + Cl\textsuperscript{-}

\[E_{\text{cell}}^\theta = 0.222 - 0.799 = -0.577\ \mathrm{V}\]

\[\lg K_{sp} = \frac{nE^\circ}{0.05916} = \frac{1 \times (-0.577)}{0.05916} = -9.75\]

\[K_{sp} = 10^{-9.75} = 1.8 \times 10^{-10}\]

\begin{center}\rule{0.5\linewidth}{0.5pt}\end{center}

\subsection{六、Latimer 图与元素电势图}\label{ux516dlatimer-ux56feux4e0eux5143ux7d20ux7535ux52bfux56fe}

\subsubsection{6.1 Latimer 图}\label{latimer-ux56fe}

按氧化态从高到低排列，标注相邻物种间的 \(E^{\circ}\)：

\[
\mathrm{MnO_4^-} \xrightarrow{+1.51} \mathrm{MnO_4^{2-}} \xrightarrow{+2.26} \mathrm{MnO_2} \xrightarrow{+0.95} \mathrm{Mn^{3+}} \xrightarrow{+1.51} \mathrm{Mn^{2+}} \xrightarrow{-1.18} \mathrm{Mn}
\]

\subsubsection{\texorpdfstring{6.2 从 Latimer 图计算非相邻物种的 \(E^{\circ}\)}{6.2 从 Latimer 图计算非相邻物种的 E\^{}\{\textbackslash circ\}}}\label{ux4ece-latimer-ux56feux8ba1ux7b97ux975eux76f8ux90bbux7269ux79cdux7684-ecirc}

对 \(\mathrm{A} \xrightarrow{E_1} \mathrm{B} \xrightarrow{E_2} \mathrm{C}\)，A→C 的 \(E^{\circ}\) 为：

\[
E^\theta_{\mathrm{A/C}} = \frac{n_1 E_1^\theta + n_2 E_2^\theta}{n_1 + n_2}
\]

\textbf{不能简单平均}---必须用电子数加权平均。这是因为 \(E^{\circ}\) 是强度量，而 \(\Delta \1^{\circ}\) 是广度量(\(\Delta \1^{\circ}\) = −nFE°)，叠加时必须按 \(\Delta \1^{\circ}\) 叠加。

\subsubsection{6.3 歧化反应的判断}\label{ux6b67ux5316ux53cdux5e94ux7684ux5224ux65ad}

在 Latimer 图中，若某物种\textbf{右侧}的 \(E^{\circ}\) \textgreater{} \textbf{左侧}的 \(E^{\circ}\)，则该物种可发生歧化反应。

\textbf{例}：Cu 的 Latimer 图(酸性介质)：

\[\mathrm{Cu^{2+}} \xrightarrow{+0.159} \mathrm{Cu^+} \xrightarrow{+0.520} \mathrm{Cu}\]

Cu\textsuperscript{+} 的歧化：2Cu\textsuperscript{+} → Cu\textsuperscript{2+} + Cu

\[E_{\text{cell}} = E_{\text{右}} - E_{\text{左}} = 0.520 - 0.159 = +0.361\ \mathrm{V} > 0\]

歧化自发！Cu\textsuperscript{+} 在酸性水溶液中\textbf{不能稳定存在}，会歧化为 Cu\textsuperscript{2+} 和 Cu。

\textbf{反例}：Fe 的 Latimer 图(酸性介质)：

\[\mathrm{Fe^{3+}} \xrightarrow{+0.77} \mathrm{Fe^{2+}} \xrightarrow{-0.44} \mathrm{Fe}\]

Fe\textsuperscript{2+} 歧化：E\_cell = E(右) − E(左) = −0.44 − 0.77 = −1.21 V \textless{} 0，\textbf{不能歧化}。Fe\textsuperscript{2+} 在酸性溶液中稳定。

\begin{center}\rule{0.5\linewidth}{0.5pt}\end{center}

\subsection{七、电化学腐蚀与防腐}\label{ux4e03ux7535ux5316ux5b66ux8150ux8680ux4e0eux9632ux8150}

\subsubsection{7.1 电化学腐蚀的类型}\label{ux7535ux5316ux5b66ux8150ux8680ux7684ux7c7bux578b}

金属腐蚀的本质是金属被氧化。电化学腐蚀需要电解质溶液和两种不同的导体(杂质或不同区域)形成微电池。

{\def\LTcaptype{none} % do not increment counter
\begin{longtable}[]{@{}llll@{}}
\toprule\noalign{}
腐蚀类型 & 氧化剂 & 酸性条件 & 反应示例 \\
\midrule\noalign{}
\endhead
\bottomrule\noalign{}
\endlastfoot
\textbf{吸氢腐蚀} & H\textsuperscript{+} & 强酸性 & Fe + 2H\textsuperscript{+} → Fe\textsuperscript{2+} + H\textsubscript{2}↑ \\
\textbf{耗氧腐蚀} & O\textsubscript{2} & 中性/弱酸性 & 2Fe + O\textsubscript{2} + 2H\textsubscript{2}O → 2Fe(OH)\textsubscript{2} \\
\end{longtable}
}

\textbf{实际中最常见的是耗氧腐蚀}，因为大气中的 O\textsubscript{2} 溶解在水膜中即可作为氧化剂。

\subsubsection{7.2 防腐方法}\label{ux9632ux8150ux65b9ux6cd5}

{\def\LTcaptype{none} % do not increment counter
\begin{longtable}[]{@{}lll@{}}
\toprule\noalign{}
方法 & 原理 & 示例 \\
\midrule\noalign{}
\endhead
\bottomrule\noalign{}
\endlastfoot
\textbf{涂层保护} & 隔绝金属与电解质 & 涂漆、涂油脂 \\
\textbf{镀层保护} & 用更耐腐蚀的金属覆盖 & 镀铬、镀锡 \\
\textbf{牺牲阳极法} & 连接更活泼的金属作为阳极 & 船体焊锌块(白铁原理) \\
\textbf{钝化处理} & 使表面形成致密氧化膜 & 不锈钢(Cr 钝化膜) \\
\end{longtable}
}

\textbf{镀层破损后的不同命运}：
- \textbf{白铁}(镀锌铁)：Zn 比 Fe 活泼，Zn 优先腐蚀保护 Fe(牺牲阳极)
- \textbf{马口铁}(镀锡铁)：Sn 比 Fe 不活泼，Sn 层破损后 Fe 加速腐蚀(Fe 成为阳极)

\begin{center}\rule{0.5\linewidth}{0.5pt}\end{center}

\subsection{八、化学电源简介}\label{ux516bux5316ux5b66ux7535ux6e90ux7b80ux4ecb}

\subsubsection{8.1 常见化学电源}\label{ux5e38ux89c1ux5316ux5b66ux7535ux6e90}

{\def\LTcaptype{none} % do not increment counter
\begin{longtable}[]{@{}
  >{\raggedright\arraybackslash}p{(\linewidth - 10\tabcolsep) * \real{0.1600}}
  >{\raggedright\arraybackslash}p{(\linewidth - 10\tabcolsep) * \real{0.1600}}
  >{\raggedright\arraybackslash}p{(\linewidth - 10\tabcolsep) * \real{0.1600}}
  >{\raggedright\arraybackslash}p{(\linewidth - 10\tabcolsep) * \real{0.1600}}
  >{\centering\arraybackslash}p{(\linewidth - 10\tabcolsep) * \real{0.2000}}
  >{\raggedright\arraybackslash}p{(\linewidth - 10\tabcolsep) * \real{0.1600}}@{}}
\toprule\noalign{}
\begin{minipage}[b]{\linewidth}\raggedright
电池类型
\end{minipage} & \begin{minipage}[b]{\linewidth}\raggedright
负极
\end{minipage} & \begin{minipage}[b]{\linewidth}\raggedright
正极
\end{minipage} & \begin{minipage}[b]{\linewidth}\raggedright
电解质
\end{minipage} & \begin{minipage}[b]{\linewidth}\centering
\(E^{\circ}\)\_cell
\end{minipage} & \begin{minipage}[b]{\linewidth}\raggedright
特点
\end{minipage} \\
\midrule\noalign{}
\endhead
\bottomrule\noalign{}
\endlastfoot
\textbf{干电池} & Zn & MnO\textsubscript{2} & NH\textsubscript{4}Cl & \textasciitilde1.5 V & 一次性，价格低 \\
\textbf{铅蓄电池} & Pb & PbO\textsubscript{2} & H\textsubscript{2}SO\textsubscript{4} & 2.04 V & 可充电，汽车常用 \\
\textbf{氢氧燃料电池} & H\textsubscript{2} & O\textsubscript{2} & KOH & 1.23 V & 清洁能源，效率高 \\
\textbf{锂离子电池} & LiC\textsubscript{6} & LiCoO\textsubscript{2} & 有机电解质 & \textasciitilde3.7 V & 高能量密度，手机/电脑 \\
\end{longtable}
}

\subsubsection{8.2 铅蓄电池详解}\label{ux94c5ux84c4ux7535ux6c60ux8be6ux89e3}

\[
\mathrm{Pb(s) + PbO_2(s) + 2H_2SO_4(aq) \rightleftharpoons 2PbSO_4(s) + 2H_2O(l)}
\]

负极：Pb(s) + SO\textsubscript{4}\textsuperscript{2-}(aq) → PbSO\textsubscript{4}(s) + 2e\textsuperscript{-}，\(E^{\circ}\) = −0.356 V
正极：PbO\textsubscript{2}(s) + 4H\textsuperscript{+}(aq) + SO\textsubscript{4}\textsuperscript{2-}(aq) + 2e\textsuperscript{-} → PbSO\textsubscript{4}(s) + 2H\textsubscript{2}O(l)，\(E^{\circ}\) = +1.685 V

\[E^\circ_{\text{cell}} = 1.685 - (-0.356) = +2.041\ \mathrm{V}\]

\[\Delta_rG^\circ = -2 \times 96485 \times 2.041 = -394\ \mathrm{kJ/mol}\]

\[\lg K^\theta = \frac{2 \times 2.041}{0.05916} = 69.0 \quad \Rightarrow \quad K^\theta \approx 10^{69}\]

\textbf{\(K^{\circ}\) 极大}，反应几乎完全进行到底。放电时 H\textsubscript{2}SO\textsubscript{4} 浓度降低，由 Nernst 方程可知 {[}H\textsuperscript{+}{]}↓ → E 正极↓ → E\_cell↓，这就是铅蓄电池放电时电压缓慢下降的原因。

\subsubsection{8.3 燃料电池}\label{ux71c3ux6599ux7535ux6c60}

以 H\textsubscript{2}-O\textsubscript{2} 碱性燃料电池为例：

负极：H\textsubscript{2} + 2OH\textsuperscript{-} → 2H\textsubscript{2}O + 2e\textsuperscript{-}
正极：O\textsubscript{2} + 2H\textsubscript{2}O + 4e\textsuperscript{-} → 4OH\textsuperscript{-}
总反应：2H\textsubscript{2} + O\textsubscript{2} → 2H\textsubscript{2}O

**燃料电池的 \(E^{\circ}\)\_cell 与 pH 无关**(因为 H\textsubscript{2} 和 O\textsubscript{2} 的半反应中 H\textsuperscript{+}/OH\textsuperscript{-} 可以互相抵消)，这是燃料电池的一大优势。

\begin{center}\rule{0.5\linewidth}{0.5pt}\end{center}

\subsection{九、典型例题}\label{ux4e5dux5178ux578bux4f8bux9898}

\textbf{题干}：写出以下电池的电极反应和电池反应，计算 \(E^{\circ}\) 和 Δ\_rG°：

\[
\text{Zn(s)} \mid \text{Zn}^{2+}\text{(1 mol/L)} \| \text{H}^+\text{(1 mol/L)} \mid \text{H}_2\text{(100 kPa)} \mid \text{Pt(s)}
\]

\textbf{解析}：

负极(氧化)：Zn(s) → Zn\textsuperscript{2+}(aq) + 2e\textsuperscript{-}，\(E^{\circ}\) = −0.763 V
正极(还原)：2H\textsuperscript{+}(aq) + 2e\textsuperscript{-} → H\textsubscript{2}(g)，\(E^{\circ}\) = 0.000 V

\[E_{\text{cell}}^\theta = 0.000 - (-0.763) = +0.763\ \mathrm{V}\]

\[\Delta_rG^\circ = -nFE^\circ = -2 \times 96485 \times 0.763 = -147.2\ \mathrm{kJ/mol}\]

\begin{center}\rule{0.5\linewidth}{0.5pt}\end{center}

\textbf{题干}：298 K 时已知：
- Ag\textsuperscript{+} + e\textsuperscript{-} ⇌ Ag(s)，\(E^{\circ}\) = +0.799 V
- Fe\textsuperscript{3+} + e\textsuperscript{-} ⇌ Fe\textsuperscript{2+}，\(E^{\circ}\) = +0.771 V

\begin{enumerate}
\def\labelenumi{(\alph{enumi})}
\tightlist
\item
  判断 Ag + Fe\textsuperscript{3+} → Ag\textsuperscript{+} + Fe\textsuperscript{2+} 是否自发。
\item
  若加入 NaCl 使 Ag\textsuperscript{+} 沉淀为 AgCl(\(K_{\1}\) = 1.8×10\textsuperscript{-10})，反应方向是否改变？
\item
  若进一步加入过量 NH\textsubscript{3} 使 Ag\textsuperscript{+} 形成 {[}Ag(NH\textsubscript{3})\textsubscript{2}{]}\textsuperscript{+}(\(K_{\1}\) = 1.1×10\textsuperscript{7})，反应方向又如何？
\end{enumerate}

\textbf{解析}：

\begin{enumerate}
\def\labelenumi{(\alph{enumi})}
\item
  E\_cell = \(E^{\circ}\)(Fe\textsuperscript{3+}/Fe\textsuperscript{2+}) − \(E^{\circ}\)(Ag\textsuperscript{+}/Ag) = 0.771 − 0.799 = −0.028 V \textless{} 0，反应\textbf{不自发}。
\item
  加入 Cl\textsuperscript{-} 后：E(AgCl/Ag) = 0.799 + 0.05916 × lg(1.8×10\textsuperscript{-10}) = +0.222 V
  新 E\_cell = 0.771 − 0.222 = +0.549 V \textgreater{} 0，反应变为\textbf{自发}！
\item
  形成 {[}Ag(NH\textsubscript{3})\textsubscript{2}{]}\textsuperscript{+} 后：游离 {[}Ag\textsuperscript{+}{]} = {[}Ag(NH\textsubscript{3})\textsubscript{2}\textsuperscript{+}{]}/(\(K_{\1}\) × {[}NH\textsubscript{3}{]}\textsuperscript{2})
  E = 0.799 + 0.05916 × lg(\(K_{\1}\)/(\(K_{\1}\) × {[}NH\textsubscript{3}{]}\textsuperscript{2})) = 0.799 − 1.052 = −0.253 V
  E\_cell = 0.771 − (−0.253) = +1.024 V \textgreater{} 0，反应\textbf{更加自发}。
\end{enumerate}

\textbf{答案}：(a) 不自发；(b) 变为自发；(c) 更加自发。

\textbf{核心洞察}：沉淀和配合物通过改变游离离子浓度来调节电极电势，从而控制反应方向。这是分析化学中''掩蔽''和''解掩蔽''的电化学基础。

\begin{center}\rule{0.5\linewidth}{0.5pt}\end{center}

\textbf{题干}：铜的 Latimer 图(酸性介质)：

\[\mathrm{Cu^{2+}} \xrightarrow{+0.159} \mathrm{Cu^+} \xrightarrow{+0.520} \mathrm{Cu}\]

\begin{enumerate}
\def\labelenumi{(\alph{enumi})}
\tightlist
\item
  计算 Cu\textsuperscript{2+}/Cu 的 \(E^{\circ}\)
\item
  判断 Cu\textsuperscript{+} 能否在酸性水溶液中稳定存在
\end{enumerate}

\textbf{解析}：

\begin{enumerate}
\def\labelenumi{(\alph{enumi})}
\tightlist
\item
  Cu\textsuperscript{2+} + 2e\textsuperscript{-} → Cu：n\textsubscript{1} = 1, E\textsubscript{1} = 0.159 V；n\textsubscript{2} = 1, E\textsubscript{2} = 0.520 V
\end{enumerate}

\[E^\circ = \frac{1 \times 0.159 + 1 \times 0.520}{2} = +0.340\ \mathrm{V}\]

\begin{enumerate}
\def\labelenumi{(\alph{enumi})}
\setcounter{enumi}{1}
\tightlist
\item
  Cu\textsuperscript{+} 的歧化：\(2\mathrm{Cu^+ \to Cu^{2+} + Cu}\)
\end{enumerate}

\[E_{\text{cell}} = E_{\text{右}} - E_{\text{左}} = 0.520 - 0.159 = +0.361\ \mathrm{V} > 0\]

歧化自发！Cu\textsuperscript{+} 在酸性水溶液中\textbf{不能稳定存在}，会歧化为 Cu\textsuperscript{2+} 和 Cu。

\begin{center}\rule{0.5\linewidth}{0.5pt}\end{center}

\textbf{题干}：铅蓄电池反应：Pb(s) + PbO\textsubscript{2}(s) + 2H\textsubscript{2}SO\textsubscript{4}(aq) ⇌ 2PbSO\textsubscript{4}(s) + 2H\textsubscript{2}O(l)
已知：\(E^{\circ}\)(PbSO\textsubscript{4}/Pb) = −0.356 V，\(E^{\circ}\)(PbO\textsubscript{2}/PbSO\textsubscript{4}) = +1.685 V。

\begin{enumerate}
\def\labelenumi{(\alph{enumi})}
\tightlist
\item
  计算 \(E^{\circ}\)\_cell 和 Δ\_rG°。
\item
  计算 \(K^{\circ}\)，解释为什么铅蓄电池在正常使用中电压基本恒定。
\item
  放电时 H\textsubscript{2}SO\textsubscript{4} 浓度降低，用 Nernst 方程解释为什么电池电压随放电而缓慢下降。
\end{enumerate}

\textbf{解析}：

\begin{enumerate}
\def\labelenumi{(\alph{enumi})}
\item
  \(E^{\circ}\)\_cell = 1.685 − (−0.356) = +2.041 V
  Δ\_rG° = −2 × 96485 × 2.041 = −393,700 J/mol ≈ −394 kJ/mol
\item
  lg \(K^{\circ}\) = 2×2.041/0.05916 = 69.0 → \(K^{\circ}\) ≈ 10\textsuperscript{69。\$K}\{\circ\}\$ 极大，反应几乎完全进行到底。由于 \(K^{\circ}\) 极大，在正常放电范围内 H\textsubscript{2}SO\textsubscript{4} 浓度变化相对较小，电压基本恒定。
\item
  放电消耗 H\textsubscript{2}SO\textsubscript{4}，{[}H\textsuperscript{+}{]} 降低。正极反应：PbO\textsubscript{2} + 4H\textsuperscript{+} + SO\textsubscript{4}\textsuperscript{2-} + 2e\textsuperscript{-} → PbSO\textsubscript{4} + 2H\textsubscript{2}O
  由 Nernst 方程，{[}H\textsuperscript{+}{]}↓ → E 正极↓ → E\_cell↓。
\end{enumerate}

\begin{center}\rule{0.5\linewidth}{0.5pt}\end{center}

\subsection{十、本节总结速查}\label{ux5341ux672cux8282ux603bux7ed3ux901fux67e5}

{\def\LTcaptype{none} % do not increment counter
\begin{longtable}[]{@{}
  >{\raggedright\arraybackslash}p{(\linewidth - 4\tabcolsep) * \real{0.3333}}
  >{\raggedright\arraybackslash}p{(\linewidth - 4\tabcolsep) * \real{0.3333}}
  >{\raggedright\arraybackslash}p{(\linewidth - 4\tabcolsep) * \real{0.3333}}@{}}
\toprule\noalign{}
\begin{minipage}[b]{\linewidth}\raggedright
核心概念
\end{minipage} & \begin{minipage}[b]{\linewidth}\raggedright
关键公式
\end{minipage} & \begin{minipage}[b]{\linewidth}\raggedright
注意点
\end{minipage} \\
\midrule\noalign{}
\endhead
\bottomrule\noalign{}
\endlastfoot
电池符号 & 负极(左) ‖ 正极(右) & 单竖线=相界面，双竖线=盐桥 \\
\(E^{\circ}\) 比较 & \(E^{\circ}\)\_cell = \(E^{\circ}\)\_正 − \(E^{\circ}\)\_负 \textgreater{} 0 自发 & \(E^{\circ}\) 是强度量，不随系数加倍 \\
Nernst 方程 & \(E = E^\theta - \frac{RT}{nF}\ln Q\) & Q 与化学平衡中的 Q 相同 \\
\(\Delta \1^{\circ}\) 与 \(E^{\circ}\) & \(\Delta_rG^\circ = -nFE^\circ_{\text{cell}}\) & n = 转移电子数 \\
\(E^{\circ}\) 与 \(K^{\circ}\) & \(E^\circ = \frac{0.05916}{n}\lg K^\theta\) & 连接热力学与电化学 \\
Latimer 图 & 氧化态从高到低排列 & 歧化判据：\(E^{\circ}\)\_右 \textgreater{} \(E^{\circ}\)\_左 \\
pH 影响 & H\textsuperscript{+} 出现在 Q 中 & pH↑ → E↓(对含 H\textsuperscript{+} 的半反应) \\
沉淀/配合物 & 降低游离离子浓度 & 使 E 向负方向偏移 \\
电化学腐蚀 & 微电池原理 & 吸氢腐蚀 vs 耗氧腐蚀 \\
\end{longtable}
}

\begin{center}\rule{0.5\linewidth}{0.5pt}\end{center}

\subsection{十一、三级练习题}\label{ux5341ux4e00ux4e09ux7ea7ux7ec3ux4e60ux9898}

\textbf{1.} 写出 Daniel 电池(Zn\textbar Zn\textsuperscript{2+}‖Cu\textsuperscript{2+}\textbar Cu)的正极反应、负极反应和电池反应。

\textbf{2.} 查表：\(E^{\circ}\)(Cu\textsuperscript{2+}/Cu) = +0.34 V，\(E^{\circ}\)(Zn\textsuperscript{2+}/Zn) = −0.76 V。计算 Daniel 电池的 \(E^{\circ}\)\_cell。

\textbf{3.} 判断下列说法的正误：
(a) \(E^{\circ}\) 是强度量，与半反应中电子数无关。
(b) \(E^{\circ}\) \textgreater{} 0 的反应在任何条件下都能自发进行。
(c) 标准氢电极的电极电势规定为 0 V。
(d) 温度升高时所有电极电势都增大。

\textbf{4.} 写出 Nernst 方程的一般形式(用 ln 表示)，说明各符号含义。

\textbf{5.} 对半反应 \(\mathrm{Fe^{3+} + e^- \rightleftharpoons Fe^{2+}}\)，\(E^{\circ}\) = +0.77 V。若 {[}Fe\textsuperscript{3+}{]} = 0.01 mol/L，{[}Fe\textsuperscript{2+}{]} = 1.0 mol/L，求 E。

\textbf{6.} 已知 \(E^{\circ}\)(Ag\textsuperscript{+}/Ag) = +0.80 V，\(E^{\circ}\)(Cu\textsuperscript{2+}/Cu) = +0.34 V。判断反应 \(\mathrm{Cu + 2Ag^+ \to Cu^{2+} + 2Ag}\) 是否自发。

\textbf{7.} 用 \(\Delta \1^{\circ}\) = −nFE° 计算第 6 题反应的 Δ\_rG°。

\textbf{8.} 简述原电池中盐桥的作用。

\textbf{9.} 什么是参比电极？列举两种常见的参比电极及其电势值。

\textbf{10.} 某电池 \(E^{\circ}\)\_cell = 0.50 V，n = 1。用 \(E^{\circ}\) 与 \(K^{\circ}\) 的关系计算 \(K^{\circ}\)(298 K)。

\begin{center}\rule{0.5\linewidth}{0.5pt}\end{center}

\textbf{11.} \textbf{(Nernst 方程综合)} 298 K 时半反应 \(\mathrm{Cu^{2+} + 2e^- \rightleftharpoons Cu}\) 的 \(E^{\circ}\) = +0.34 V。
(a) 求 {[}Cu\textsuperscript{2+}{]} = 0.001 mol/L 时的 E。
(b) 若向 Cu\textsuperscript{2+} 溶液中加入过量 NH\textsubscript{3} 使 {[}Cu(NH\textsubscript{3})\textsubscript{4}\textsuperscript{2+}{]} = 0.1 mol/L，{[}NH\textsubscript{3}{]} = 1.0 mol/L，已知 \(K_{\1}\) = 1.1×10\textsuperscript{13}，求新的 E。

\textbf{12.} \textbf{(沉淀对电极电势的影响)} 已知 \(E^{\circ}\)(Pb\textsuperscript{2+}/Pb) = −0.126 V，\(K_{\1}\)(PbSO\textsubscript{4}) = 1.6×10\textsuperscript{-8}。计算 \(E^{\circ}\)(PbSO\textsubscript{4}/Pb)。解释为什么铅蓄电池放电时 PbSO\textsubscript{4} 沉淀的生成有利于维持电压。

\textbf{13.} \textbf{(\(\Delta \1^{\circ}\) − \(E^{\circ}\) − \(K^{\circ}\) 三角)} 某氧化还原反应的 \(E^{\circ}\)\_cell = 0.30 V，n = 2。分别计算 Δ\_rG° 和 \(K^{\circ}\)(298 K)。

\textbf{14.} \textbf{(Latimer 图分析)} 铁在酸性介质中的 Latimer 图：

\[\mathrm{Fe^{3+}} \xrightarrow{+0.77} \mathrm{Fe^{2+}} \xrightarrow{-0.44} \mathrm{Fe}\]

\begin{enumerate}
\def\labelenumi{(\alph{enumi})}
\tightlist
\item
  计算 Fe\textsuperscript{3+}/Fe 的 \(E^{\circ}\)。
\item
  判断 Fe\textsuperscript{2+} 能否发生歧化。
\item
  计算 Fe + 2Fe\textsuperscript{3+} → 3Fe\textsuperscript{2+} 的 \(E^{\circ}\)\_cell 和 \(K^{\circ}\)。
\end{enumerate}

\textbf{15.} \textbf{(pH 对电极电势的影响)} 半反应 \(\mathrm{Cr_2O_7^{2-} + 14H^+ + 6e^- \rightleftharpoons 2Cr^{3+} + 7H_2O}\)，\(E^{\circ}\) = +1.33 V。
(a) 写出该半反应的 Nernst 方程。
(b) 计算 pH = 7 时的 E(假设其他物种为标准浓度)。
(c) 解释为什么重铬酸钾在碱性溶液中氧化能力显著减弱。

\textbf{16.} \textbf{(电池设计)} 设计一个原电池来计算 \(K_{\1}\)(BaSO\textsubscript{4})，写出电池符号、电极反应，并推导 \(K_{\1}\) 与 \(E^{\circ}\) 的关系。已知 \(E^{\circ}\)(Ba\textsuperscript{2+}/Ba) = −2.90 V，\(E^{\circ}\)(SO\textsubscript{4}\textsuperscript{2-}/SO\textsubscript{2}) = +0.17 V。

\textbf{17.} \textbf{(浓差电池)} 构建浓差电池 Cu\textbar Cu\textsuperscript{2+}(0.001 mol/L)\textbar\textbar Cu\textsuperscript{2+}(1.0 mol/L)\textbar Cu。计算 E\_cell(298 K)，并解释为什么浓差电池的电动势随两侧浓度趋近而减小。

\textbf{18.} \textbf{(电化学与热力学统一)} 已知 298 K 时：
- \(\mathrm{Fe^{2+} + 2e^- \rightleftharpoons Fe}\)，\(E^{\circ}\) = −0.44 V
- \(\mathrm{O_2 + 4H^+ + 4e^- \rightleftharpoons 2H_2O}\)，\(E^{\circ}\) = +1.23 V

\begin{enumerate}
\def\labelenumi{(\alph{enumi})}
\tightlist
\item
  写出 Fe 在酸性溶液中被 O\textsubscript{2} 氧化的电池反应。
\item
  计算 \(E^{\circ}\)\_cell 和 Δ\_rG°。
\item
  计算 \(K^{\circ}\)，解释为什么铁在酸性含氧水中会缓慢腐蚀。
\end{enumerate}

\textbf{19.} \textbf{(Nernst 与 pH 图应用)} 对于反应 \(\mathrm{Zn + 2H^+ \to Zn^{2+} + H_2}\)，\(E^{\circ}\)(Zn\textsuperscript{2+}/Zn) = −0.76 V，\(E^{\circ}\)(H\textsuperscript{+}/H\textsubscript{2}) = 0 V。
(a) 计算标准状态下反应的 \(E^{\circ}\)\_cell。
(b) 若 {[}Zn\textsuperscript{2+}{]} = 1.0 mol/L，pH = 3，p(H\textsubscript{2}) = 100 kPa，计算 E\_cell。
(c) 若用非氧化性酸(如 HCl)，反应实际能否发生？讨论动力学因素的影响。

\begin{center}\rule{0.5\linewidth}{0.5pt}\end{center}

\textbf{20.} \textbf{(全国初赛)} 298 K 时已知：
- \(\mathrm{MnO_4^- + 8H^+ + 5e^- \rightleftharpoons Mn^{2+} + 4H_2O}\)，\(E^{\circ}\) = +1.51 V
- \(\mathrm{MnO_4^- + 2H_2O + 3e^- \rightleftharpoons MnO_2 + 4OH^-}\)，\(E^{\circ}\) = +0.59 V
- \(K_w = 1.0 \times 10^{-14}\)

\begin{enumerate}
\def\labelenumi{(\alph{enumi})}
\tightlist
\item
  计算 MnO\textsubscript{2}/Mn\textsuperscript{2+} 在酸性介质中的 \(E^{\circ}\)(即 \(\mathrm{MnO_2 + 4H^+ + 2e^- \rightleftharpoons Mn^{2+} + 2H_2O}\) 的 \(E^{\circ}\))。
\item
  计算 MnO\textsubscript{2} 在 pH = 14 的碱性溶液中歧化的 \(E^{\circ}\)\_cell 和 \(K^{\circ}\)。
\item
  解释为什么 KMnO\textsubscript{4} 在强碱性溶液中还原产物是 MnO\textsubscript{2} 而非 Mn\textsuperscript{2+}。
\end{enumerate}

\textbf{21.} \textbf{(全国初赛改编)} 铅蓄电池反应：\(\mathrm{Pb(s) + PbO_2(s) + 2H_2SO_4(aq) \rightleftharpoons 2PbSO_4(s) + 2H_2O(l)}\)。
已知：\(E^{\circ}\)(PbSO\textsubscript{4}/Pb) = −0.356 V，\(E^{\circ}\)(PbO\textsubscript{2}/PbSO\textsubscript{4}) = +1.685 V。
(a) 计算 \(E^{\circ}\)\_cell。
(b) 计算 Δ\_rG°。
(c) 计算 \(K^{\circ}\)，解释为什么铅蓄电池在正常使用中电压基本恒定。
(d) 放电时 H\textsubscript{2}SO\textsubscript{4} 浓度降低，用 Nernst 方程解释为什么电池电压随放电而缓慢下降。

\textbf{22.} \textbf{(综合计算)} 298 K 时：
- \(\mathrm{Ag^+ + e^- \rightleftharpoons Ag}\)，\(E^{\circ}\) = +0.799 V
- \(\mathrm{Fe^{3+} + e^- \rightleftharpoons Fe^{2+}}\)，\(E^{\circ}\) = +0.771 V
- \(\mathrm{Fe^{2+} + 2e^- \rightleftharpoons Fe}\)，\(E^{\circ}\) = −0.440 V

\begin{enumerate}
\def\labelenumi{(\alph{enumi})}
\tightlist
\item
  判断 \(\mathrm{Ag + Fe^{3+} \to Ag^+ + Fe^{2+}}\) 是否自发。
\item
  若向体系中加入 NaCl 使 Ag\textsuperscript{+} 沉淀为 AgCl(\(K_{\1}\) = 1.8×10\textsuperscript{-10})，反应方向是否改变？
\item
  若进一步加入过量 NH\textsubscript{3} 使 Ag\textsuperscript{+} 形成 {[}Ag(NH\textsubscript{3})\textsubscript{2}{]}\textsuperscript{+}(\(K_{\1}\) = 1.1×10\textsuperscript{7})，反应方向又如何？
\end{enumerate}

\textbf{23.} \textbf{(电化学与溶解平衡)} 已知 \(E^{\circ}\)(Cu\textsuperscript{2+}/Cu) = +0.34 V，\(E^{\circ}\)(Cu\textsubscript{2}O/Cu) = −0.36 V，\(K_{\1}\) = 1.0×10\textsuperscript{-14}。
(a) 计算 Cu\textsubscript{2}O 在酸性溶液中歧化的 \(E^{\circ}\)\_cell(\(\mathrm{Cu_2O + 2H^+ \to Cu + Cu^{2+} + H_2O}\))。
(b) 计算 Cu\textsubscript{2}O 在纯水中的溶解平衡常数。
(c) 讨论 Cu\textsubscript{2}O 在酸性和碱性水溶液中的稳定性差异。

\textbf{24.} \textbf{(标准电极电势的热力学本质)} 证明：对于半反应 \(\mathrm{Ox + ne^- \rightleftharpoons Red}\)，标准电极电势可表示为：

\[E^\theta = \frac{\Delta_f G^\theta(\mathrm{Ox}) - \Delta_f G^\theta(\mathrm{Red})}{nF}\]

并用此公式从 Δ\_fG° 数据计算 \(E^{\circ}\)(Cu\textsuperscript{2+}/Cu)。

\textbf{25.} \textbf{(综合设计题)} 你需要通过电化学方法测定 BaSO\textsubscript{4} 的 \(K_{\1}\)。
(a) 设计一个原电池(写出电池符号和电极反应)。
(b) 推导 \(E^{\circ}\)\_cell 与 \(K_{\1}\) 的关系式。
(c) 若测得 \(E^{\circ}\)\_cell = 3.25 V，计算 \(K_{\1}\)(BaSO\textsubscript{4})。
(d) 讨论该方法的实验可行性(需要考虑哪些实际因素？)。

\begin{center}\rule{0.5\linewidth}{0.5pt}\end{center}

\subsection{十二、练习题答案}\label{ux5341ux4e8cux7ec3ux4e60ux9898ux7b54ux6848}

\textbf{1.}
负极：Zn(s) → Zn\textsuperscript{2+}(aq) + 2e\textsuperscript{-}
正极：Cu\textsuperscript{2+}(aq) + 2e\textsuperscript{-} → Cu(s)
电池反应：Zn(s) + Cu\textsuperscript{2+}(aq) → Zn\textsuperscript{2+}(aq) + Cu(s)

\textbf{2.} \(E^{\circ}\)\_cell = 0.34 − (−0.76) = +1.10 V

\textbf{3.}
(a) \(E^{\circ}\) 是强度量，与电子数无关。
(b) 错误。\(E^{\circ}\) \textgreater{} 0 只在标准状态下判断。非标态需用 Nernst 方程。
(c) SHE 电势规定为 0 V。
(d) 错误。温度对 E 的影响取决于具体半反应，不是所有都增大。

\textbf{4.} \(E = E^\theta + \frac{RT}{nF}\ln\frac{[\mathrm{Ox}]}{[\mathrm{Red}]}\)，其中 \(E^{\circ}\) 为标准电极电势，R = 8.314 J/(mol·K)，T 为温度(K)，n 为转移电子数，F = 96485 C/mol，{[}Ox{]}/{[}Red{]} 为氧化态/还原态的相对浓度。

\textbf{5.} E = 0.77 + 0.05916×lg(0.01/1.0) = 0.77 − 0.118 = +0.652 V

\textbf{6.} \(E^{\circ}\)\_cell = \(E^{\circ}\)(Ag\textsuperscript{+}/Ag) − \(E^{\circ}\)(Cu\textsuperscript{2+}/Cu) = 0.80 − 0.34 = +0.46 V \textgreater{} 0，自发。

\textbf{7.} Δ\_rG° = −nFE° = −2 × 96485 × 0.46 = −88,767 J/mol ≈ −88.8 kJ/mol

\textbf{8.} 盐桥维持两半电池的电荷平衡。没有盐桥，负极侧因 Zn\textsuperscript{2+} 积累而带正电，正极侧因 Cu\textsuperscript{2+} 消耗而带负电，电荷积累会阻碍电子流动，电池很快停止工作。

\textbf{9.} 参比电极是电极电势已知且稳定的电极，用于测量其他电极的电势。常见：标准氢电极(SHE，0 V)、Ag/AgCl 电极(+0.222 V)、饱和甘汞电极(SCE，+0.244 V)。

\textbf{10.} lg \(K^{\circ}\) = nE°/0.05916 = 1×0.50/0.05916 = 8.45。\(K^{\circ}\) = 10\^{}8.45 ≈ 2.8 × 10\textsuperscript{8}。

\begin{center}\rule{0.5\linewidth}{0.5pt}\end{center}

\textbf{11.}
(a) E = 0.34 + (0.05916/2)×lg(0.001) = 0.34 − 0.089 = +0.251 V

\begin{enumerate}
\def\labelenumi{(\alph{enumi})}
\setcounter{enumi}{1}
\tightlist
\item
  游离 {[}Cu\textsuperscript{2+}{]} = {[}Cu(NH\textsubscript{3})\textsubscript{4}\textsuperscript{2+}{]}/(\(K_{\1}\) × {[}NH\textsubscript{3}{]}\textsuperscript{4}) = 0.1/(1.1×10\textsuperscript{13} × 1.0) = 9.1×10\textsuperscript{-15} mol/L
  E = 0.34 + (0.05916/2)×lg(9.1×10\textsuperscript{-15}) = 0.34 − 0.413 = −0.073 V
\end{enumerate}

配合物形成使 Cu\textsuperscript{2+}/Cu 的电势从 +0.34 V 降至 −0.07 V。

\textbf{12.} \(E^{\circ}\)(PbSO\textsubscript{4}/Pb) = \(E^{\circ}\)(Pb\textsuperscript{2+}/Pb) + (0.05916/2)×lg(\(K_{\1}\)) = −0.126 + (0.05916/2)×lg(1.6×10\textsuperscript{-8}) = −0.126 − 0.230 = −0.356 V

放电时 Pb 被氧化为 PbSO\textsubscript{4}(沉淀)，使 {[}Pb\textsuperscript{2+}{]} 降低。Nernst 方程预测 E(Pb\textsuperscript{2+}/Pb) 增大(向正方向)，有利于维持电池电压。

\textbf{13.} Δ\_rG° = −2 × 96485 × 0.30 = −57,891 J/mol ≈ −57.9 kJ/mol
\(K^{\circ}\) = e\^{}(2×96485×0.30/(8.314×298)) = e\^{}23.4 ≈ 1.4 × 10\textsuperscript{10}

\textbf{14.}
(a) \(E^{\circ}\)(Fe\textsuperscript{3+}/Fe) = (1×0.77 + 2×(−0.44))/3 = (0.77−0.88)/3 = −0.037 V

\begin{enumerate}
\def\labelenumi{(\alph{enumi})}
\setcounter{enumi}{1}
\item
  Fe\textsuperscript{2+} 歧化：E\_cell = E(右) − E(左) = −0.44 − 0.77 = −1.21 V \textless{} 0，\textbf{不能歧化}。Fe\textsuperscript{2+} 在酸性溶液中稳定。
\item
  E\_cell = \(E^{\circ}\)(Fe\textsuperscript{3+}/Fe\textsuperscript{2+}) − \(E^{\circ}\)(Fe\textsuperscript{2+}/Fe) = 0.77 − (−0.44) = +1.21 V
  \(K^{\circ}\) = e\^{}(2×96485×1.21/(8.314×298)) ≈ e\^{}94.3 ≈ 10\^{}41
\end{enumerate}

\textbf{15.}
(a) \(E = 1.33 + \frac{0.05916}{6}\lg\frac{[\mathrm{Cr_2O_7^{2-}}][\mathrm{H^+}]^{14}}{[\mathrm{Cr^{3+}}]^2}\)

\begin{enumerate}
\def\labelenumi{(\alph{enumi})}
\setcounter{enumi}{1}
\item
  pH = 7 时 {[}H\textsuperscript{+}{]} = 10\textsuperscript{-7}：E = 1.33 + (0.05916/6)×lg(10\textsuperscript{-7})\textsuperscript{14} = 1.33 − 0.966 = +0.364 V
\item
  碱性溶液中 H\textsuperscript{+} 浓度极低，Cr\textsubscript{2}O\textsubscript{7}\textsuperscript{2-} 转化为 CrO\textsubscript{4}\textsuperscript{2-}，且 {[}H\textsuperscript{+}{]}\^{}\{14\} 项使 E 大幅降低，氧化能力显著减弱。
\end{enumerate}

\textbf{16.} 设计电池：Ba(s)\textbar Ba\textsuperscript{2+}(1 mol/L)\textbar\textbar SO\textsubscript{4}\textsuperscript{2-}(1 mol/L)\textbar PbSO\textsubscript{4}(s)\textbar Pb(s)

更实际的设计：Pt\textbar H\textsubscript{2}(100 kPa)\textbar H\textsuperscript{+}(1 mol/L)\textbar\textbar Ba\textsuperscript{2+}(1 mol/L)\textbar BaSO\textsubscript{4}(s)\textbar Ba(s)

需要已知 \(E^{\circ}\)(SO\textsubscript{4}\textsuperscript{2-}/SO\textsubscript{2}) = +0.17 V。通过两个半反应组合求 \(K_{\1}\)。

实际上更简单：使用 Ag/AgCl 类似方法。

电池：Ag\textbar Ag\textsuperscript{+}(1 mol/L)\textbar\textbar Ba\textsuperscript{2+}(饱和 BaSO\textsubscript{4})\textbar BaSO\textsubscript{4}(s)\textbar Ba(s)

\textbf{17.} E\_cell = (0.05916/2)×lg(1.0/0.001) = (0.05916/2)×3 = +0.089 V

当两侧浓度趋近时，Q → 1，lg Q → 0，E\_cell → 0。浓差电池的驱动力就是浓度梯度。

\textbf{18.}
(a) 电池反应：2Fe(s) + O\textsubscript{2}(g) + 4H\textsuperscript{+}(aq) → 2Fe\textsuperscript{2+}(aq) + 2H\textsubscript{2}O(l)
(b) \(E^{\circ}\)\_cell = 1.23 − (−0.44) = +1.67 V
Δ\_rG° = −4 × 96485 × 1.67 = −644,390 J/mol ≈ −644 kJ/mol
(c) lg \(K^{\circ}\) = 4×1.67/0.05916 = 113 → \(K^{\circ}\) ≈ 10\^{}113，极大。铁在含氧酸性水中腐蚀在热力学上极其有利，但动力学上腐蚀速率取决于氧的扩散速率。

\textbf{19.}
(a) \(E^{\circ}\)\_cell = 0 − (−0.76) = +0.76 V
(b) E = 0.76 + (0.05916/2)×lg(1.0×(10\textsuperscript{-3})\textsuperscript{2}) = 0.76 − 0.178 = +0.58 V
(c) 热力学上 E\_cell \textgreater{} 0，反应自发。但锌与非氧化性酸反应速率取决于氢气在锌表面的析出动力学(氢过电位较高)，实际反应可能较慢。锌越纯，反应越慢(杂质形成微电池加速腐蚀)。

\begin{center}\rule{0.5\linewidth}{0.5pt}\end{center}

\textbf{20.}
(a) 方法：将两个已知半反应组合求目标 \(E^{\circ}\)。
MnO\textsubscript{4}\textsuperscript{-} + 8H\textsuperscript{+} + 5e\textsuperscript{-} → Mn\textsuperscript{2+} + 4H\textsubscript{2}O \ldots{} E\textsubscript{1} = +1.51 V, n\textsubscript{1} = 5
MnO\textsubscript{4}\textsuperscript{-} + 2H\textsubscript{2}O + 3e\textsuperscript{-} → MnO\textsubscript{2} + 4OH\textsuperscript{-} \ldots{} E\textsubscript{2} = +0.59 V, n\textsubscript{2} = 3

先将(2)在碱性条件下转化为酸性：MnO\textsubscript{4}\textsuperscript{-} + 4H\textsuperscript{+} + 3e\textsuperscript{-} → MnO\textsubscript{2} + 2H\textsubscript{2}O，E\textsubscript{2} = ?

由 E\textsubscript{2}(碱) = +0.59 V，使用 Nernst 方程处理 {[}OH\textsuperscript{-}{]} 与 {[}H\textsuperscript{+}{]} 的关系：
在 pH = 0 时 E = 0.59 + (0.05916/3)×lg({[}OH\textsuperscript{-}{]}\textsuperscript{4}) = 0.59 + (0.05916/3)×lg(10\textsuperscript{-56})

更直接的方法：用 ΔG 叠加。
ΔG\textsubscript{1} = −5F × 1.51 = −7.55F
ΔG\textsubscript{2} = −3F × 0.59 = −1.77F(碱性条件下)
转化为酸性需要修正 OH\textsuperscript{-}/H\textsuperscript{+} 的影响。

简化处理：目标反应 = (1) − (2) 适当调整。
MnO\textsubscript{2} + 4H\textsuperscript{+} + 2e\textsuperscript{-} → Mn\textsuperscript{2+} + 2H\textsubscript{2}O

\(E^{\circ}\) = (n\textsubscript{1}E\textsubscript{1} − n\textsubscript{2}E\textsubscript{2})/(n\textsubscript{1} − n\textsubscript{2}) = (5×1.51 − 3×0.59)/(5−3)\ldots{} 这不正确因为(2)是碱性条件。

实际上在酸性条件下：MnO\textsubscript{4}\textsuperscript{-} + 4H\textsuperscript{+} + 3e\textsuperscript{-} → MnO\textsubscript{2} + 2H\textsubscript{2}O，\(E^{\circ}\) = +1.70 V(文献值)
目标 \(E^{\circ}\)(MnO\textsubscript{2}/Mn\textsuperscript{2+}) = (5×1.51 − 3×1.70)/2 = (7.55 − 5.10)/2 = +1.225 V

\begin{enumerate}
\def\labelenumi{(\alph{enumi})}
\setcounter{enumi}{1}
\item
  MnO\textsubscript{2} 在碱中歧化：2MnO\textsubscript{2} + 2OH\textsuperscript{-} → MnO\textsubscript{4}\textsuperscript{-} + MnO\textsubscript{2}\textsuperscript{2-} + H\textsubscript{2}O(不常见)
  更典型的是 MnO\textsubscript{2}\textsuperscript{2-} 歧化。
\item
  KMnO\textsubscript{4} 在强碱中还原产物为 MnO\textsubscript{2}(+4 价)，因为碱性条件下 MnO\textsubscript{2}/Mn\textsuperscript{2+} 的 \(E^{\circ}\) 降低，MnO\textsubscript{4}\textsuperscript{-}/MnO\textsubscript{2} 的 \(E^{\circ}\) 仍较高，MnO\textsubscript{4}\textsuperscript{-} 最稳定的还原产物变为 MnO\textsubscript{2}。
\end{enumerate}

\textbf{21.}
(a) \(E^{\circ}\)\_cell = 1.685 − (−0.356) = +2.041 V
(b) Δ\_rG° = −2 × 96485 × 2.041 = −393,700 J/mol ≈ −394 kJ/mol
(c) lg \(K^{\circ}\) = 2×2.041/0.05916 = 69.0 → \(K^{\circ}\) ≈ 10\textsuperscript{69。\$K}\{\circ\}\$ 极大，反应几乎完全进行到底。
(d) 放电消耗 H\textsubscript{2}SO\textsubscript{4}，{[}H\textsuperscript{+}{]} 降低。由 Nernst 方程，H\textsuperscript{+} 出现在正极反应(PbO\textsubscript{2} + 4H\textsuperscript{+} + SO\textsubscript{4}\textsuperscript{2-} + 2e\textsuperscript{-} → PbSO\textsubscript{4} + 2H\textsubscript{2}O)中，{[}H\textsuperscript{+}{]}↓ → E 正极↓ → E\_cell↓。

\textbf{22.}
(a) E\_cell = \(E^{\circ}\)(Fe\textsuperscript{3+}/Fe\textsuperscript{2+}) − \(E^{\circ}\)(Ag\textsuperscript{+}/Ag) = 0.771 − 0.799 = −0.028 V \textless{} 0，反应\textbf{不自发}。

\begin{enumerate}
\def\labelenumi{(\alph{enumi})}
\setcounter{enumi}{1}
\item
  加入 Cl\textsuperscript{-} 后 Ag\textsuperscript{+} 沉淀：E(AgCl/Ag) = +0.222 V
  新 E\_cell = 0.771 − 0.222 = +0.549 V \textgreater{} 0，反应变为\textbf{自发}！
\item
  形成 {[}Ag(NH\textsubscript{3})\textsubscript{2}{]}\textsuperscript{+} 后：E = 0.799 + 0.05916×lg(\(K_{\1}\)/(\(K_{\1}\)×{[}NH\textsubscript{3}{]}\textsuperscript{2}))
  = 0.799 + 0.05916×(−17.8) = 0.799 − 1.052 = −0.253 V
  E\_cell = 0.771 − (−0.253) = +1.024 V \textgreater{} 0，反应\textbf{更加自发}。
\end{enumerate}

\textbf{23.}
(a) Cu\textsubscript{2}O + 2H\textsuperscript{+} → Cu + Cu\textsuperscript{2+} + H\textsubscript{2}O
E\_cell = \(E^{\circ}\)(Cu\textsuperscript{2+}/Cu\textsubscript{2}O) − \(E^{\circ}\)(Cu\textsuperscript{+}/Cu\textsubscript{2}O)

由 Cu\textsubscript{2}O/Cu: Cu\textsubscript{2}O + 2H\textsuperscript{+} + 2e\textsuperscript{-} → 2Cu + H\textsubscript{2}O，E = −0.36 V(已知)
由 Cu\textsubscript{2}O/Cu\textsuperscript{2+}: Cu\textsubscript{2}O + 2H\textsuperscript{+} → 2Cu\textsuperscript{2+} + H\textsubscript{2}O + 2e\textsuperscript{-}(逆向)

需要用 Cu\textsuperscript{2+}/Cu\textsubscript{2}O 和 Cu\textsuperscript{+}/Cu\textsubscript{2}O 的数据。Cu\textsubscript{2}O 歧化的 \(E^{\circ}\)\_cell ≈ +0.46 V(文献值)。

\begin{enumerate}
\def\labelenumi{(\alph{enumi})}
\setcounter{enumi}{1}
\item
  Cu\textsubscript{2}O 溶解：Cu\textsubscript{2}O + H\textsubscript{2}O ⇌ 2Cu\textsuperscript{+} + 2OH\textsuperscript{-}
  K = {[}Cu\textsuperscript{+}{]}\textsuperscript{2}{[}OH\textsuperscript{-}{]}\textsuperscript{2}(极小，Cu\textsubscript{2}O 在水中几乎不溶)
\item
  酸性中 Cu\textsubscript{2}O 溶解(H\textsuperscript{+} 与 OH\textsuperscript{-} 中和驱动)，碱性中 Cu\textsubscript{2}O 稳定(溶解平衡向左)。
\end{enumerate}

\textbf{24.} 对半反应 Ox + ne\textsuperscript{-} → Red：
由热力学：Δ\_rG° = Δ\_fG°(Red) − Δ\_fG°(Ox)(标准反应的生成焓差)
又 Δ\_rG° = −nFE°

\[E^\theta = \frac{\Delta_f G^\theta(\mathrm{Ox}) - \Delta_f G^\theta(\mathrm{Red})}{nF}\]

注意符号：Δ\_rG° = −nFE°，所以 \(E^{\circ}\) = −Δ\_rG°/(nF) = {[}Δ\_fG°(Ox) − Δ\_fG°(Red){]}/(nF)

对 Cu\textsuperscript{2+} + 2e\textsuperscript{-} → Cu：Δ\_fG°(Cu) = 0，Δ\_fG°(Cu\textsuperscript{2+}) = +64.98 kJ/mol
\(E^{\circ}\) = (64980 − 0)/(2×96485) = +0.337 V ≈ +0.34 V ✓

\textbf{25.}
(a) 电池设计：需要两个半电池分别涉及 BaSO\textsubscript{4} 的溶解平衡和一个已知参考电极。
方案：Pt\textbar H\textsubscript{2}(100 kPa)\textbar H\textsuperscript{+}(1 mol/L)\textbar\textbar Ba\textsuperscript{2+}(饱和 BaSO\textsubscript{4})\textbar BaSO\textsubscript{4}(s)\textbar Ba(s)

\begin{enumerate}
\def\labelenumi{(\alph{enumi})}
\setcounter{enumi}{1}
\item
  \(E^{\circ}\)\_cell 与 \(K_{\1}\) 的关系需要通过多个已知 \(E^{\circ}\) 推导。设 \(E^{\circ}\)(Ba\textsuperscript{2+}/Ba) = −2.90 V，目标是求 \(K_{\1}\)。
  更实际的方案：使用类似 Ag/AgCl 的方法。
\item
  需要具体数据才能计算。
\item
  实际困难：Ba 非常活泼(\(E^{\circ}\) 极负)，会与水剧烈反应；需要在无水或非水介质中操作；参比电极的盐桥可能引入干扰离子。
\end{enumerate}

\end{document}
